\color{unidarkgreen}{\subsection{Логика сигнализации}}
\color{black}

\begin{enumerate}

\item
Устройство формирует выходное воздействие на сигнализацию \textit{«Сигнализация: Срабатывание фикс.»} и \textit{«Сигнализация: Срабатывание»} по схеме «ИЛИ» в следующих случаях:
\begin{itemize}
\item при срабатывании КИВ;
\item при срабатывании дифференциальных защит;
\item при срабатывании газовых и технологических защит;
\item при срабатывании общей логики отключения; 
\item при срабатывании ЗНР;
\item при срабатывании УРОВ.
\end{itemize}

\item
Устройство формирует выходное воздействие на сигнализацию \textit{«Сигнализация: Неисправность фикс.»} и \textit{«Сигнализация: Неисправность»} по схеме «ИЛИ» в следующих случаях:
\begin{itemize}
\item при неисправности КИВ;
\item при неисправности токовых цепей дифференциальных защит;
\item при неисправности газовых и технологических защит;
\item при неисправности оперативного тока;
\item при пуске ЗПО;
\item при срабатывании датчиков технологической сигнализации (уровень масла, неисправность охлаждения и пр.);
\item при неисправности выключателя.
\end{itemize}


\item
Алгоритм работы логики сигнализации приведен на рисунке \ref{pict_ls:LS}.


\medskip
\begin{figure}[H]
\centering
\includegraphics[width=0.7\textwidth,height=0.7\textheight,keepaspectratio]{\fbpath/11.01 ШР/91 Сигнализация/img/ls.pdf}
\captionof{figure}{Функциональная схема логики сигнализации}\label{pict_ls:LS}
\end{figure}

\end{enumerate}